% auto-generated by experiment_combined_ablation.py
\begin{table*}[t]
\centering
\small
\caption{Critical runtime dependencies}
\label{tab:table_combined_ablation_C}
\begin{tabular}{lll}
\toprule
Component & Dependent Components & Failure Impact (measured single removal) \\
\midrule
Predicate Engine (PE) & EQ+FT+RD & RIS$\to$0.83; FPR 1.000; evidence 1.000; ledger 1.000 \\
Runtime Revocation (RV) & FT+RD & RIS$\to$0.83; FPR 0.519; evidence 1.000; ledger 1.000 \\
Evidence Quad (EQ) & HC+LG & RIS$\to$0.33; FPR 0.519; evidence 0.000; ledger n/a \\
Runtime Ledger (LG) & HC & RIS$\to$0.50; FPR 0.519; evidence 1.000; ledger n/a \\
Hash Chain (HC) & — & RIS$\to$0.50; FPR 0.519; evidence 1.000; ledger 0.000 \\
Runtime Risk Detection (RD) & — & governance stage (not in ablation matrix) \\
Runtime Watchdog (WD) & FT & governance stage (not in ablation matrix) \\
Fleet Telemetry (FT) & — & governance stage (not in ablation matrix) \\
Clock Consistency (single-host PTP) (CK) & — & governance stage (not in ablation matrix) \\
\bottomrule
\end{tabular}
\\[3pt]
\begin{minipage}{\textwidth}\footnotesize \textbf{Note --- URR is not the False Permit Rate.} $\mathrm{URR}\;(\text{Undetected Risk Rate}) = \mathrm{FN}/(\mathrm{TP}+\mathrm{FN}) = 1 - \text{Blind Detection Recall}$. This metric measures the fraction of malicious events that remain undetected during blind runtime evaluation. It is \emph{not} the False Permit Rate reported in the main authorization benchmark. The paper's False Permit Rate ($0/492$ and $0/62$) measures authorization soundness and remains unchanged.\end{minipage}
\end{table*}
